\documentclass[8pt, a4paper]{article}

% ── Packages ─────────────────────────────────────────────────────────────────
\usepackage[a4paper, top=2.5cm, bottom=2.5cm, left=2.8cm, right=2.8cm]{geometry}
\usepackage[T1]{fontenc}
\usepackage[utf8]{inputenc}
\usepackage{lmodern}
\usepackage{microtype}
\usepackage{parskip}          % space between paragraphs, no indent
\usepackage{titlesec}         % section heading styles
\usepackage{enumitem}         % list customisation
\usepackage{booktabs}         % professional tables
\usepackage{tabularx}         % auto-width table columns
\usepackage{array}            % extra column formatting
\usepackage{xcolor}           % colours
\usepackage{mdframed}         % framed boxes
\usepackage{fancyhdr}         % header / footer
\usepackage{hyperref}         % clickable links (load last)

% ── Colours ───────────────────────────────────────────────────────────────────
\definecolor{accentblue}{RGB}{30, 90, 160}
\definecolor{rulegray}{RGB}{190, 195, 205}
\definecolor{boxbg}{RGB}{237, 242, 250}

% ── Hyperref setup ────────────────────────────────────────────────────────────
\hypersetup{
  colorlinks = true,
  linkcolor  = accentblue,
  urlcolor   = accentblue,
}

% ── Section heading styles ────────────────────────────────────────────────────
\titleformat{\section}
  {\large\bfseries\color{accentblue}}
  {\thesection.}{0.6em}{}
  [\vspace{-4pt}\textcolor{rulegray}{\rule{\linewidth}{0.6pt}}\vspace{2pt}]

\titlespacing{\section}{0pt}{18pt}{6pt}

% ── Concept box ───────────────────────────────────────────────────────────────
\newmdenv[
  backgroundcolor   = boxbg,
  linecolor         = accentblue,
  linewidth         = 1.2pt,
  leftline          = true,
  rightline         = false,
  topline           = false,
  bottomline        = false,
  innerleftmargin   = 10pt,
  innertopmargin    = 6pt,
  innerbottommargin = 6pt,
  skipabove         = 6pt,
  skipbelow         = 6pt,
]{conceptbox}

% ── List styling ──────────────────────────────────────────────────────────────
\setlist[itemize]{topsep=3pt, itemsep=2pt, parsep=0pt, leftmargin=1.6em}
\setlist[enumerate]{topsep=3pt, itemsep=2pt, parsep=0pt, leftmargin=1.8em}

% ── Header / Footer ───────────────────────────────────────────────────────────
\pagestyle{fancy}
\fancyhf{}
\renewcommand{\headrulewidth}{0.4pt}
\fancyhead[L]{\small\color{accentblue}\textbf{ELEC3609 / ELEC9609}}
\fancyhead[R]{\small\color{gray}Week 1 --- Introduction \& Course Mechanics}
\fancyfoot[C]{\small\thepage}

% =============================================================================
\begin{document}

% ── Title Block ───────────────────────────────────────────────────────────────
\begin{center}
  {\LARGE\bfseries\color{accentblue} ELEC3609 / ELEC9609: Internet Software Platforms}\\[6pt]
  {\large Week 1 Study Guide: Introduction \& Course Mechanics}
\end{center}

\vspace{12pt}
\textcolor{rulegray}{\rule{\linewidth}{1pt}}
\vspace{6pt}

% =============================================================================
\section{Course Overview \& Mechanics}

\textbf{Course Code:} ELEC3609 (Undergraduate) / ELEC9609 (Postgraduate)\\
\textbf{Course Title:} Internet Software Platforms\\
\textbf{Course Coordinator \& Lecturer:} Dr.\ Huaming Chen (School of ECE)\\
\textbf{Primary Tech Stack:} Python, Django Framework, HTML/CSS/JavaScript, REST APIs, AWS Cloud

\medskip

\begin{conceptbox}
\textbf{Key Learning Outcomes (LOs)}
\begin{itemize}
  \item \textbf{LO1:} Use Python \& Django as baseline programming tools.
  \item \textbf{LO2--LO4:} Design, build, test, and maintain full-stack web-based applications from inception to market.
  \item \textbf{LO5:} Deploy web apps using cloud services (AWS), DevOps practices, and baseline security measures.
  \item \textbf{LO6--LO9:} Practice team collaboration, technical writing, ethical responsibilities, accessibility, and cybersecurity.
\end{itemize}
\end{conceptbox}

\medskip

\begin{conceptbox}
\textbf{Assessment Breakdown}

\smallskip
\begin{tabularx}{\linewidth}{X c l}
  \toprule
  \textbf{Assessment Task} & \textbf{Weight} & \textbf{Due Timeline} \\
  \midrule
  Weekly Individual Tasks           & 8\%  & Weeks 2--4, Weeks 6--10 \\
  System \& Database Design         & 15\% & W5 (Presentation) \& W6 (Report) \\
  System Implementation \& Testing  & 35\% & Week 11 \\
  System Deployment \& Security     & 12\% & Week 13 \\
  Supervised Final Exam             & 30\% & Final Exam Period \\
  \bottomrule
\end{tabularx}
\end{conceptbox}

\medskip

\textbf{Group Project Workflow} (Groups of up to 5 students):

\smallskip\noindent
Idea Concept $\to$ Requirements \& Use Cases $\to$ UI Wireframes $\to$ Architecture \& Database Design $\to$ Coding \& Testing $\to$ Cloud Deployment $\to$ Final Project Presentation

% =============================================================================
\section{Core Concepts: The Internet vs.\ the World Wide Web (WWW)}

\begin{conceptbox}
\textbf{Concept 1: Internet}
\begin{itemize}
  \item \textbf{Definition:} A global network of networks connecting millions of computing devices.
  \item \textbf{Purpose:} Provides the physical and logical medium allowing computers to communicate with each other.
\end{itemize}
\end{conceptbox}

\begin{conceptbox}
\textbf{Concept 2: World Wide Web (WWW or ``The Web'')}
\begin{itemize}
  \item \textbf{Definition:} An information-sharing system built ON TOP of the Internet.
  \item \textbf{Access Method:} Information is accessed via documents/pages connected by web links.
  \item \textbf{Crucial Distinction:} The Web is NOT the Internet! It is just ONE of many applications/services operating over the Internet.
  \item Other non-Web Internet services include:
  \begin{itemize}
    \item Email (SMTP / IMAP / POP3)
    \item Instant Messaging (WebSocket / XMPP)
    \item File Transfer (FTP / SFTP)
    \item Media Streaming (RTSP / RTMP)
    \item Usenet / SSH / Peer-to-Peer (P2P)
  \end{itemize}
\end{itemize}
\end{conceptbox}

\begin{conceptbox}
\textbf{Concept 3: Websites, URLs, and URIs}
\begin{itemize}
  \item \textbf{Website:} A collection of related web pages hosted on the Web, constructed using:
  \begin{itemize}
    \item \textbf{HTML} (HyperText Markup Language): Defines structural content.
    \item \textbf{CSS} (Cascading Style Sheets): Defines visual styling.
    \item \textbf{JavaScript:} Defines client-side interactivity and logic.
  \end{itemize}
  \item \textbf{URI vs.\ URL:}
  \begin{itemize}
    \item \textbf{URI} (Uniform Resource Identifier): A generic term for all names and addresses that identify a resource.
    \item \textbf{URL} (Uniform Resource Locator): A specific TYPE of URI that provides the location and method of accessing a resource (e.g., protocol + domain + path).
    \item \textbf{Rule of Thumb:} All URLs are URIs, but not all URIs are URLs!
  \end{itemize}
  \item \textbf{Standard URL Structure:}\\
    \texttt{protocol://domain.tld/path/file.ext}\\
    \textit{Example:} \texttt{https://google.com} (HTTPS protocol), \texttt{mailto:user@sydney.edu.au}, \texttt{ftp://ftp.google.com}
\end{itemize}
\end{conceptbox}

% =============================================================================
\section{Internet Software Platforms: The ``Stitching'' Concept}

An Internet Software Platform stitches together three foundational pillars:

\begin{enumerate}
  \item \textbf{The Internet (Infrastructure \& Communication):}
  \begin{itemize}
    \item Networking fundamentals \& data transport
    \item Protocols (HTTP, TCP/IP, REST)
    \item Cybersecurity fundamentals (Encryption, Auth, TLS)
  \end{itemize}

  \item \textbf{Software (Engineering \& Architecture):}
  \begin{itemize}
    \item Design patterns and architectural principles
    \item Coding (Front-end \& Back-end, e.g., Django/Python)
    \item Packaging \& Modularization
    \item Human factors (UX/UI, usability, accessibility)
  \end{itemize}

  \item \textbf{Platforms (Cloud \& Operations):}
  \begin{itemize}
    \item Cloud technologies (IaaS, PaaS, SaaS)
    \item Infrastructure as Code (IaC --- managing infrastructure via scripts)
    \item DevOps \& DevSecOps (Automated pipelines, continuous integration, security integration)
  \end{itemize}
\end{enumerate}

% =============================================================================
\section{Software Platform Lifecycle \& DevOps in Practice}

\begin{conceptbox}
\textbf{Software Platform Lifecycle}

\medskip
The life cycle of an internet software platform alternates continuously between two main phases:

\medskip
\textbf{A. Development Phase:}
\begin{itemize}
  \item Plan \& Design (Architecture, DB schemas, wireframes)
  \item Build \& Test (Coding, Unit testing, Integration testing)
  \item Package (Containerization e.g., Docker, building artifacts)
\end{itemize}

\textbf{B. Operations Phase:}
\begin{itemize}
  \item Release \& Configure (Deployment to servers/cloud)
  \item Monitor (Tracking performance, uptime, latency, logs)
  \item React \& Respond (Patching bugs, scaling resources, incident response)
\end{itemize}
\end{conceptbox}

\begin{conceptbox}
\textbf{Real-World Data Pipelines (ETL Architecture)}

\smallskip
Modern web platforms handle massive streams of data through pipelines:
\begin{enumerate}
  \item \textbf{Data Ingestion:} Sensor data or user events ingested in real-time, near real-time, or batch.
  \item \textbf{Data Processing:} Extract, Transform, and Load (ETL) processing.
  \item \textbf{Data Analytics:} Performing computations or machine learning inference.
  \item \textbf{Data Visualization:} Displaying insights to users on web dashboards.
\end{enumerate}
\end{conceptbox}

\begin{conceptbox}
\textbf{Industry Spotlight: DevOps \& Automation (e.g., Qantas Case Study)}
\begin{itemize}
  \item \textbf{Scalability with Automation:} A small platform team (e.g., 12 engineers) can manage 250+ applications and millions of server instances using automation.
  \item \textbf{Everything as Code:} Infrastructure, security, and deployment pipelines are written as code scripts rather than manual setup.
  \item \textbf{Blue/Green Deployments:} Technique for zero-downtime deployment by running two identical environments (Blue = Live production, Green = Idle/New version). Traffic switches seamlessly to Green once verified.
  \item \textbf{Site Reliability Engineering (SRE):} Applying software engineering practices to infrastructure and operations problems.
\end{itemize}
\end{conceptbox}

% =============================================================================
\section{Key Criteria When Building Websites/Platforms}

Taking an idea from conception to market requires balancing 10 major factors:

\begin{enumerate}
  \item \textbf{Cost:} Hosting, API usage, computational resources.
  \item \textbf{Performance:} Load times, latency, throughput.
  \item \textbf{Hosting:} Cloud providers (AWS, GCP, Azure), serverless vs.\ dedicated.
  \item \textbf{Stability / Reliability:} Uptime, failover mechanisms, redundancy.
  \item \textbf{Payment:} Payment gateway integration (Stripe, PayPal).
  \item \textbf{Stack / Frameworks:} Choosing tools like Django, React, Postgres, Redis.
  \item \textbf{Maintenance:} Code quality, modular design, documentation.
  \item \textbf{Security:} Data protection, HTTPS, vulnerability prevention (OWASP Top 10).
  \item \textbf{Audience:} User scale, internationalization, accessibility.
  \item \textbf{Compliance / Ethics:} Privacy laws (GDPR), data governance.
\end{enumerate}

% =============================================================================
\section{Quick Summary \& Memory Cheat Sheet}

\begin{itemize}
  \item \textbf{Internet vs.\ Web:} Internet = Physical \& Network Highway; Web = One vehicle running on that highway (HTTP web pages).
  \item \textbf{URL vs.\ URI:} URI = Identifier (Name/ID); URL = Locator (Address + Protocol).
  \item \textbf{The 3 Pillars of ELEC3609:} Internet (Networking/Security) + Software (Django/Design) + Platform (AWS Cloud/DevOps).
  \item \textbf{Core Lifecycle:} Plan $\to$ Build $\to$ Package $\to$ Release $\to$ Monitor $\to$ React.
  \item \textbf{Deployment Pattern:} Blue/Green Deployment = Zero downtime deployment between 2 identical environments.
  \item \textbf{ETL:} Extract, Transform, Load (Data Ingestion pipeline).
\end{itemize}

% =============================================================================
\section{Self-Assessment Questions for Retention}

\textbf{Q1: Explain the fundamental difference between the Internet and the World Wide Web.}

\smallskip
\textit{Answer:} The Internet is the global network infrastructure connecting machines; the World Wide Web is an application/service running on top of the Internet using HTTP to share hyperlinked web pages.

\bigskip

\textbf{Q2: What is the relationship between URIs and URLs?}

\smallskip
\textit{Answer:} Every URL is a URI, but not every URI is a URL. A URI identifies a resource; a URL identifies it and describes how to access it (location + protocol).

\bigskip

\textbf{Q3: What are the two major phases of an Internet Software Platform's lifecycle?}

\smallskip
\textit{Answer:} Development (Plan, Build/Test, Package) and Operations (Release/Configure, Monitor, React/Respond).

\bigskip

\textbf{Q4: What is Blue/Green deployment?}

\smallskip
\textit{Answer:} A deployment strategy utilizing two identical physical/virtual environments. One serves active traffic (Blue) while the other receives the new version (Green). After testing, live traffic is switched to Green, enabling zero-downtime releases and instant rollbacks.

\vspace{14pt}
\textcolor{rulegray}{\rule{\linewidth}{1pt}}

\end{document}
